\documentclass[11pt,letterpaper]{article}

\usepackage[utf8]{inputenc}
\usepackage[T1]{fontenc}
\usepackage{lmodern}
\usepackage{microtype}
\usepackage[margin=1in]{geometry}
\usepackage{xcolor}
\usepackage{titlesec}
\usepackage{enumitem}
\usepackage{listings}
\usepackage{tcolorbox}
\usepackage{hyperref}
\usepackage{booktabs}

% Palette definition
\definecolor{primaryblue}{RGB}{30, 58, 138}
\definecolor{slateborder}{RGB}{203, 213, 225}
\definecolor{codebg}{RGB}{248, 250, 252}
\definecolor{promptbg}{RGB}{241, 245, 249}
\definecolor{accentdark}{RGB}{15, 23, 42}
\definecolor{keywordblue}{RGB}{2, 132, 199}
\definecolor{commentgreen}{RGB}{22, 101, 52}

\tcbuselibrary{skins,breakable}

% Hyperref setup
\hypersetup{
    colorlinks=true,
    linkcolor=primaryblue,
    urlcolor=keywordblue,
    citecolor=primaryblue
}

% Section styling
\titleformat{\section}
  {\normalfont\Large\bfseries\color{primaryblue}}
  {\thesection}{1em}{}[\titlerule]

\titleformat{\subsection}
  {\normalfont\large\bfseries\color{accentdark}}
  {\thesubsection}{1em}{}

% Listings configuration for Bash/Git
\lstdefinestyle{gitbash}{
    backgroundcolor=\color{codebg},
    basicstyle=\small\ttfamily\color{accentdark},
    breaklines=true,
    breakatwhitespace=true,
    frame=single,
    rulecolor=\color{slateborder},
    framesep=6pt,
    tabsize=2,
    showstringspaces=false,
    commentstyle=\color{commentgreen}\itshape,
    keywordstyle=\color{keywordblue}\bfseries,
    morekeywords={git, checkout, branch, stash, reset, clean, status, agy, function, return}
}
\lstset{style=gitbash}

% Custom tcolorbox for Prompts
\newtcolorbox{promptbox}[1][]{
    colback=promptbg,
    colframe=slateborder,
    coltitle=accentdark,
    fonttitle=\bfseries\sffamily,
    title={Original Query},
    sharp corners,
    boxrule=0.8pt,
    left=12pt,
    right=12pt,
    top=8pt,
    bottom=8pt,
    breakable,
    #1
}

% Custom tcolorbox for Recovery Alert/Callout
\newtcolorbox{recoverybox}[1][]{
    colback=codebg,
    colframe=primaryblue,
    coltitle=white,
    fonttitle=\bfseries\sffamily,
    title={Recovery Hatch Summary},
    sharp corners,
    boxrule=1pt,
    left=10pt,
    right=10pt,
    top=8pt,
    bottom=8pt,
    #1
}

\begin{document}

\begin{center}
    {\LARGE \textbf{\color{primaryblue}Git State Management and Recovery Discipline}}\\[0.4em]
    {\large \textbf{Resilient Operational Patterns for Google Antigravity 2.0 CLI (\texttt{agy}) Sessions}}\\[0.8em]
    \rule{\linewidth}{0.8pt}
\end{center}

\vspace{0.5em}

\section*{Context and User Inquiry}

\begin{promptbox}
``If my Ubuntu 24.04 VM were unstable I would consider using tmux, but it has not crashed since LLMs helped me make it stable. So I will not likely need tmux. I only worry about leaving the repo in unknown state if the system reboots during agy sessions that are saving work to the repo. Is there a git discipline that makes it easier to return to known good state after a disconnect or something else that leaves the repo in an unknown state? Can we versions or release numbers or tags or something? I am not very familiar with git let alone a discipline that might be useful here.''
\end{promptbox}

\section{Introduction: Branches vs. Tags and Version Numbers}

When managing autonomous code-generation workflows using agentic command-line tools like \texttt{agy}, unexpected interruptions—such as an operating system reboot, network drop, or virtual machine shutdown—can occur in the middle of multi-file write operations. This runs the risk of leaving the repository in a corrupted, inconsistent, or untracked state.

A common intuition is to apply semantic versioning numbers or Git tags before running an agent. However:
\begin{itemize}[leftmargin=1.5em, itemsep=0.3em]
    \item \textbf{Git Tags and Version Numbers} represent immutable milestones intended to mark permanent releases (e.g., \texttt{v1.0.0}, \texttt{v1.2.3-rc1}) in project history. They are metadata pointers created to look backward at completed, verified work.
    \item \textbf{Ephemeral Feature Branches} provide dynamic, disposable sandboxes designed specifically for in-flight, exploratory, or automated mutations.
\end{itemize}

Using a dedicated branching discipline rather than version tags eliminates the risk of an unknown repository state, enabling an instantaneous rollback to a known good state with a single command.

\section{The Three-Pillar Git Discipline}

\subsection{Pillar 1: Always Isolate Agent Work in a Disposable Branch}
Never invoke \texttt{agy} directly on a trunk branch such as \texttt{main} or \texttt{master}. Prior to starting an agent session, branch off into a dedicated sandbox prefixed with a clear namespace (e.g., \texttt{agy/}):

\begin{lstlisting}[language=bash]
# Create and switch to a disposable agent branch
git checkout -b agy/session-task-name
\end{lstlisting}

If an abrupt reboot occurs mid-session, the integrity of \texttt{main} remains intact. If the session output is determined to be malformed or unrecoverable, the entire branch can be discarded:

\begin{lstlisting}[language=bash]
# Abandon the corrupted session and return to safety
git checkout main
git branch -D agy/session-task-name
\end{lstlisting}

\subsection{Pillar 2: Verify a Pristine Starting State}
An autonomous agent that starts in a repository containing dirty unstaged changes or loose files will interweave new modifications with existing work. Before running \texttt{agy}, ensure that the working tree is clean:

\begin{lstlisting}[language=bash]
# Check working directory state
git status

# If uncommitted edits exist, stow them away cleanly
git stash -u   # -u includes untracked files
\end{lstlisting}

\subsection{Pillar 3: The Deterministic One-Command Recovery Hatch}
If the system restarts unexpectedly and you return to find half-written files, partial patches, or syntax errors, there is no need to manually audit what the agent touched. 

Because the session started from a known clean commit on the branch, the repository can be returned to that exact state instantly:

\begin{lstlisting}[language=bash]
# 1. Revert all modified tracked files to the initial commit
git reset --hard HEAD

# 2. Remove all newly created untracked files and directories
git clean -fd
\end{lstlisting}

\begin{recoverybox}
\textbf{The Golden Guarantee:} Running \texttt{git reset --hard HEAD \&\& git clean -fd} completely purges every uncommitted file alteration and newly created artifact, resetting the workspace to the exact byte-level state of the last recorded commit.
\end{recoverybox}

\section{Standard Operating Procedure (SOP) Lifecycle}

The recommended operational cycle for running \texttt{agy} sessions without state contamination follows a structured four-stage pattern:

\begin{enumerate}[leftmargin=2em, itemsep=0.8em]
    \item \textbf{Pre-Flight (Baseline):} Ensure your base branch is up to date and clean.
    \begin{lstlisting}[language=bash]
    git checkout main
    git pull origin main
    git status  # Must show 'working tree clean'
    \end{lstlisting}

    \item \textbf{Branch \& Launch:} Isolate the execution context.
    \begin{lstlisting}[language=bash]
    git checkout -b agy/refactor-pipeline
    agy
    \end{lstlisting}

    \item \textbf{Post-Session Triage (Post-Flight):}
    \begin{itemize}[leftmargin=1.5em, itemsep=0.3em]
        \item \textbf{Scenario A: Interrupted / Corrupted.} Run the recovery hatch:
        \begin{lstlisting}[language=bash]
        git reset --hard HEAD
        git clean -fd
        \end{lstlisting}
        \item \textbf{Scenario B: Successful Run.} Inspect diffs, test, and commit:
        \begin{lstlisting}[language=bash]
        git diff
        git add -A
        git commit -m "feat(pipeline): automated refactor via agy"
        \end{lstlisting}
    \end{itemize}

    \item \textbf{Integration or Teardown:}
    \begin{itemize}[leftmargin=1.5em, itemsep=0.3em]
        \item If accepted: Merge back into \texttt{main} via fast-forward or squash merge:
        \begin{lstlisting}[language=bash]
        git checkout main
        git merge agy/refactor-pipeline
        git branch -d agy/refactor-pipeline
        \end{lstlisting}
        \item If rejected: Force delete the branch without touching \texttt{main}:
        \begin{lstlisting}[language=bash]
        git checkout main
        git branch -D agy/refactor-pipeline
        \end{lstlisting}
    \end{itemize}
\end{enumerate}

\section{Automating the Workflow with a BASH Helper}

To avoid repetitive manual branch creation and pre-flight checks, add the following helper function to your \texttt{\textasciitilde/.bashrc}:

\begin{lstlisting}[language=bash]
# Safe wrapper for Google Antigravity CLI sessions
agy_sandbox() {
    local task_name="${1:-scratch}"
    local branch_name="agy/${task_name}-$(date +%Y%m%d-%H%M%S)"

    # Ensure repository is clean
    if ! git diff-index --quiet HEAD -- 2>/dev/null; then
        echo "Error: Working directory has uncommitted changes. Stash or commit first."
        return 1
    fi

    echo "Creating sandbox branch: ${branch_name}"
    git checkout -b "${branch_name}" || return 1

    echo "Launching Antigravity CLI..."
    agy

    echo ""
    echo "Session completed or exited."
    echo "Current branch: ${branch_name}"
    echo "To keep work:    git add -A && git commit -m '...'"
    echo "To reset state:  git reset --hard HEAD && git clean -fd"
    echo "To abandon:      git checkout main && git branch -D ${branch_name}"
}
\end{lstlisting}

\section{Summary Reference Table}

\begin{center}
\begin{tabular}{@{}lll@{}}
\toprule
\textbf{Goal} & \textbf{Git Command} & \textbf{Operational Effect} \\ \midrule
Isolate agent session & \texttt{git checkout -b agy/<task>} & Prevents pollution of \texttt{main} \\
Stash local work & \texttt{git stash -u} & Preserves prior uncommitted work \\
Audit agent changes & \texttt{git status} / \texttt{git diff} & Reviews modified and new files \\
Revert tracked edits & \texttt{git reset --hard HEAD} & Restores tracked files to base \\
Purge new artifacts & \texttt{git clean -fd} & Deletes unversioned files/dirs \\
Abandon entire attempt & \texttt{git checkout main \&\& git branch -D <branch>} & Leaves zero trace of session \\ \bottomrule
\end{tabular}
\end{center}

\end{document}
